% SSSUP-FONTE-STORICA: fonti-locali/originali/2014-2020.zip
% Anno 2020, esercizio 2. Testo storico; figura ricostruita da PDF.
% Verificare editorialmente la trasposizione prima di considerarla definitiva.
\renewcommand{\Anno}{2020}
\renewcommand{\NumeroProblema}{2}

\renewcommand{\Testo}{
Siano $a,b$ due numeri scelti a caso nell'intervallo $[0,1]$.
\begin{enumerate}
\item[(a)] Determinare la probabilità $p_3$ che il polinomio
$$P_3(x)=x^3-ax+b$$
abbia tre radici reali distinte.
\item[(b)] Per ogni intero $n\ge3$ determinare la probabilità $p_n$ che il polinomio
$$P_n(x)=x^n-ax+b$$
abbia tre radici reali distinte.
\item[(c)] Calcolare il limite delle probabilità
$$\lim_{n\to\infty}p_{2n+1}.$$
\end{enumerate}
}

\renewcommand{\Risposta}{}
\renewcommand{\Aiutino}{}
\renewcommand{\SoluzioneUfficiale}{}
\renewcommand{\SoluzioneRedazionale}{}
\renewcommand{\SoluzioniAlternative}{}
\renewcommand{\Osservazioni}{}
\renewcommand{\Commenti}{}
